\section{From record scales to all time scales}
\label{tree:sec:bootstrap}

The preceding comparison and rigidity arguments apply to subsequences
with suitable control of their normalization.  We now prove that this
control holds at every sufficiently large scale.  The proof has three
steps: compare the internal endpoint norm with the metric multiplier
in a rigid limit; fix strict tests that hold near every such limit;
and rule out a first scale where those tests fail.

Throughout this section, an internal interval cost uses only the
instructions witnessed in that interval.  It need not equal the
unrestricted distance between the tips of the limiting cell.  Recall
\[
 X_I=d_I(\text{first tip},\text{last tip}),\qquad
 A_j=\|X_{[0,2^j]}\|_k,\qquad s=p/2.
\]
The normalizations and the constants attached to a probability law
are deterministic.  Stationarity and concatenation give
\begin{equation}
 0<A_j<\infty,\qquad A_{j+1}\leq 2A_j .
 \label{tree:eq:bootstrap_elementary_growth}
\end{equation}
Every extraction retains the complete interval and certificate data,
including marked subpaths.  This preserves the internal normalization
while the ambient passage rule is being identified.

\subsection{Comparing the internal norm with the metric multiplier}

Call a certificate law \emph{rigid} if its local passages agree
almost surely with $C D_h$ for one deterministic $C\in(0,\infty)$,
as in Theorem~\ref{tree:thm:rigidity}.  Our first task is to compare
this multiplier with the internal endpoint norm used to normalize
the discrete model.
For such a law let $q_i$ be the $L^k$ norm of its internal endpoint
cost on $[0,2^{-i}]$, in its base cost units.  These are the actual
limiting moments supplied by Lemma~\ref{tree:lem:moment_ui}, not merely
upper bounds from Fatou's lemma.  We consider rigid laws for which
\begin{equation}
 H:=\sup_{i\geq0}2^{ip}q_i<\infty.
 \label{tree:eq:bootstrap_H}
\end{equation}
Both alternatives of Proposition~\ref{tree:prop:comparison_families},
followed by Theorem~\ref{tree:thm:rigidity}, give laws of this kind.

\begin{proposition}[Universal comparison of internal norms]
\label{tree:prop:internal_norm_comparison}
There are constants $0<\alpha\leq\beta<\infty$, depending only on
the fixed Brownian and metric units and on $k$, with the following
property.  For every rigid certificate law satisfying
\eqref{tree:eq:bootstrap_H}, and every fixed integer $i$, its internal
endpoint norm satisfies
\begin{equation}
 \alpha C2^{-ip}\leq q_i\leq\beta C2^{-ip}.
 \label{tree:eq:bootstrap_norm_bridge}
\end{equation}
Negative $i$ mean a finite enlargement of the time interval.  The
assertion applies also to laws obtained by joint extraction of zooms
of already limiting laws.  Consequently, if $q_0=1$, then
\begin{equation}
 C_-:=\beta^{-1}\leq C\leq\alpha^{-1}=:C_+.
 \label{tree:eq:bootstrap_C_range}
\end{equation}
\end{proposition}

\begin{proof}
\smallskip\noindent\emph{A lower bound from the two tips.}
Let $W=D_h(\eta(0),\eta(1))$ in the fixed unit cone law.  Its tips
are distinct almost surely, so $W>0$ almost surely.  Choose
$a>0$ and $b>0$ with $\Prob[W\geq a]\geq b$.  The lower comparison
for certificates, applied in successively larger open neighborhoods
of the unit cell, gives $X_{[0,1]}\geq CW$.  Hence
\begin{equation}
 q_0\geq C a b^{1/k}.
 \label{tree:eq:bootstrap_lower_tip_norm}
\end{equation}
This argument requires no integrability statement about $W$ at a
cell tip.

\smallskip\noindent\emph{An upper bound by a maximizing zoom.}
It suffices to bound $H/C$ over the laws under consideration.
If no such bound exists, choose laws indexed by $r$ with
$H_r/C_r\to\infty$.
For each $r$ choose an integer $i_r\geq0$ such that
\[
 2^{i_rp}q_{r,i_r}>H_r/2,
 \qquad v_r=2^{-i_r}.
\]
Change time units to $v_r$ and cost units to $H_rv_r^p$.  This
is the quantum-unit transport of Proposition~\ref{tree:prop:zoom_transport};
its kernel is the transported law, and need not be the original
kernel.  In the new units the internal endpoint norms $\widetilde
q_{r,l}$ obey
\begin{equation}
 \tfrac12<\widetilde q_{r,0}\leq1,
 \qquad \widetilde q_{r,l}\leq2^{-lp}\quad(l\geq0).
 \label{tree:eq:bootstrap_maximizing_zoom}
\end{equation}
Stationarity and addition extend these estimates, with constants
depending only on the enlargement, to each bounded time window.
Lemmas~\ref{tree:lem:moment_compactness} and~\ref{tree:lem:moment_ui}, and
Proposition~\ref{tree:prop:certificate_limit}, therefore give a further
joint extraction with a limiting unit endpoint variable
$\widetilde X$ satisfying
\[
 \E[\widetilde X^k]\geq2^{-k}.
\]
The uniform integrability argument remains applicable here: internal
costs in disjoint intervals have the product interval law, and their
mixed monomials are controlled exactly as in
Lemma~\ref{tree:lem:moment_ui}.  Indeed every finite vector of costs
on deterministic disjoint intervals has an unconditional product
law at discrete resolution, and hence after each fixed zoom
extraction.  Products of weakly convergent measures converge to
the product, so the same assertion survives a sequence of different
continuum laws and maximizing zooms.  Thus this extraction does
not lose the positive endpoint moment; no closure of conditional
independence is being inferred in this moment argument.

In the same units, the ambient upper comparison is
$(C_r/H_r)D_h$.  The underlying continuum data have the fixed cone
law after quantum-unit transport.  The upper realization statement
with open margins in Proposition~\ref{tree:prop:certificate_limit} sends
this comparison to zero on every compact local route.  More
explicitly, work first on a countable exhaustion by buffered charts
and with bounded reference metric costs.  Finite upper certificates
realize the route in each chart; their costs tend to zero.  Passing
to the closed certificate sets and then removing the reference-cost
cutoff makes every such limiting local cost zero almost surely.
This is a joint statement with the reference metric; it does not
use continuity of the metric as a function of the contour.

Lemma~\ref{tree:lem:nonzero_local_cost} now gives a contradiction.  Its
geometric content is useful here.  Two points at curve times
approaching $0$ and $1$ can be chosen in the interior of the unit
disk cell.  A finite chain of compact interior routes connects them.
The routes lift using the cell's internal instructions, so their
internal cost is zero in the limit.  The two omitted tip pieces
have arbitrarily small cost by the uniform time modulus.  The
unit internal endpoint cost must therefore vanish, contrary to
$\E[\widetilde X^k]\geq2^{-k}$.  We have proved $H\leq\beta C$
for some finite universal $\beta$.  Together with
\eqref{tree:eq:bootstrap_lower_tip_norm}, and decreasing $\alpha$ if
necessary, this proves \eqref{tree:eq:bootstrap_norm_bridge} for $i=0$.

\smallskip\noindent\emph{Every fixed dyadic interval.}
For any fixed dyadic $v>0$, change time units to $v$ and divide
costs by $v^p$.  The reference metric has again the fixed cone
units and the comparison constant is again $C$.  The transported
law satisfies \eqref{tree:eq:bootstrap_H}.  For $v\leq1$ this follows
directly from that equation.  For $v>1$, all but finitely many of
its dyadic norms are controlled by that equation, while the
remaining norms are finite by addition on a bounded enlargement.
Applying the already proved unit-interval statement to this law
proves \eqref{tree:eq:bootstrap_norm_bridge} for every integer $i$.
This also proves the assertion about finite enlargements.  At no
point was a Weyl identity for the competing law assumed.
\end{proof}

\subsection{A fixed neighborhood of the rigid laws}

We now choose tests that every rigid limit passes with a strict
margin.  The resulting neighborhood and numerical thresholds will
remain fixed throughout the first-failure argument.

\paragraph{Compactness of the rigid class.}
Let $\mathcal R$ be the closure, in the topology of
Section~\ref{tree:sec:coherent-background}, of the rigid
unit-internal-norm laws satisfying \eqref{tree:eq:bootstrap_H},
including their admissible joint extractions.  This is a compact
set.  To verify the assertion, Proposition~\ref{tree:prop:internal_norm_comparison}
gives the uniform bounds
\begin{equation}
 q_0=1,\qquad q_i\leq(\beta/\alpha)2^{-ip}
 \quad(i\geq0),\qquad C\in[C_-,C_+].
 \label{tree:eq:bootstrap_rigid_compact}
\end{equation}
The moment and certificate compactness lemmas give precompactness
on every bounded enlargement.  Uniform integrability preserves
$q_0=1$ and the bounds in \eqref{tree:eq:bootstrap_rigid_compact}.
The retained reference-law coordinates and
Lemma~\ref{tree:lem:certificate_class_closed} preserve the admissible
local kernels through this further extraction; the full
background is attached by
\eqref{tree:eq:bootstrap_background_closure}.
After a further extraction, the deterministic constants $C$
converge within $[C_-,C_+]$.  Upper realization and lower
subextraction with open margins preserve the two comparisons
with this limiting constant.  The limit is therefore again
rigid.  This proves compactness and also explains why
$\mathcal R$ retains all internal and certificate data: it is
not obtained by discarding those data and keeping only the
one-parameter family of ambient metrics.

\paragraph{Numerical thresholds and the ambient compact class.}
Put $R_0=\beta/\alpha$.  Choose $R$ and then an integer
$m\geq1$ satisfying
\begin{equation}
 R>2\max\{2,2^pR_0,2^{-p}R_0\},\qquad
 R_0 2^{-m(p-s)}<\tfrac12.
 \label{tree:eq:bootstrap_constants}
\end{equation}
Set
\[
 K_0=\max\{R_0,1\}(R2^s)^{m-1}.
\]
The number $K_0$ is fixed before any induction or stopping
argument.  It bounds the past tail along a successful run, as
shown below.  For choosing neighborhoods one can use a larger
fixed constant, say $K_1=4(1+2^s+2^p)K_0$, to include one
finite forward enlargement in preceding units.

Here is the compactness fact needed to choose the neighborhood.
Let $\mathcal P_{K_1}$ be the closure of the continuum tree
certificate laws with the fixed background marginal, the limiting
interval bounds of Lemma~\ref{tree:lem:moment_compactness} with
constant $K_1$, and the retained reference-law extension. The
internal interval coordinates carry the compatibility inherited
from the walk: their laws are stationary under deterministic time
translation, costs concatenate, and costs on every fixed finite
list of deterministic intervals with disjoint interiors have
their product law. These
properties pass through joint weak convergence of the retained
endpoint coordinates, as in the maximizing-zoom argument above.
The contour coordinate agrees with the retained Brownian
background contour.

The interval bounds imply uniform integrability within this class
by the proof of Lemma~\ref{tree:lem:moment_ui}. Indeed, subdivide a
fixed rational interval into $\ell$ equal pieces. In the multinomial
expansion of the sum of their costs to the $k$th power, the pure
terms have total expectation at most
$C_TK_1^k \ell^{1-ks}$. For each fixed $\ell$, independence bounds every
mixed term uniformly in $L^{k/(k-1)}$. First send the tail threshold
to infinity and then $\ell$ to infinity. This proves the required
uniform integrability since $ks>1$.

Moment bounds on the countable windows and the fixed tight
background marginal give precompactness. The uniform
integrability just proved, Lemma~\ref{tree:lem:certificate_class_closed},
and \eqref{tree:eq:bootstrap_background_closure} preserve the
certificate and kernel properties in the closure. Thus
$\mathcal P_{K_1}$ is compact. It includes the limiting laws of
eventually successful runs: each fixed past offset has the
$K_0$ bound, even if a finite initial portion of a discrete
sequence has a larger constant. The set $\mathcal R$ is contained
in this compact class. A finite-resolution $\mu_j$ lives in the
ambient law space and need not belong to $\mathcal P_{K_1}$.

\paragraph{Stable shell tests.}
Choose the shell geometry and the high probability threshold
required by the second alternative of
Proposition~\ref{tree:prop:comparison_families}.  There are fixed
finite upper connection and positive lower traversal constants
for which every law in $\mathcal R$ passes those shell tests
with strictly stronger constants and a strictly smaller failure
probability.  In fact its local rule is $C D_h$ with
$C\in[C_-,C_+]$; compactness, continuity and nondegeneracy of
$D_h$ on buffered shells give the choices, uniformly in $C$.
Include the reference boundary-diameter and modulus tests used
in that proposition.  These choices involve only compact
interiors of the shells.

For completeness, the passage from this statement to a weak-law
neighborhood uses both halves of certificate convergence.
All shell predicates here are predicates on continuum laws
in $\mathcal P_{K_1}$.  Finite-resolution laws will be tested
only for membership in the ambient open neighborhood below.
An upper connection bound with strict cost and spatial slack
is realized by finitely many certificates for a finite net;
the uniform joining modulus extends it to every point of the
band.  A violation of a lower traversal bound gives a
bounded-cost certificate with separated marked endpoints;
compact subextraction preserves the violation at a slightly
weaker threshold.  Increasing finite nets and decreasing
joining errors gives these statements at all points.  The
uniform modulus in $\mathcal P_{K_1}$ permits the nets and
errors to be chosen before passing to the limit.  In using
the time modulus on a spatial compact $K\Subset U$, first
enlarge $K$ slightly inside $U$.  Properness of $\eta$ makes
the preimage of this enlargement lie in a finite random
time window.  The common background marginal permits a
deterministic window to be chosen with arbitrarily small
failure probability, uniformly over $\mathcal P_{K_1}$.
On that window use the interval modulus and the compact
joining argument of Proposition~\ref{tree:prop:certificate_limit},
then remove the properness cutoff.  Truncated
density approximation gives the same implications for the
specified local free-field tests.  Consequently, any sequence
in $\mathcal P_{K_1}$ converging to a law in $\mathcal R$
eventually satisfies the weaker shell tests.
These are the uniform upper and marked lower implications of
Lemma~\ref{tree:lem:certificate_buffered_tests}.  Lower tests do
not require a bound on all possible crossings: a violation
is witnessed by one bounded-cost path with its crossing
endpoints marked before extraction.

Let $\mathcal F$ be the closure in $\mathcal P_{K_1}$ of the
laws failing these weaker tests.  The preceding sequential
statement, with the intervening strict slack, gives
$\mathcal F\cap\mathcal R=\varnothing$.  Both sets are compact.
If $\mathcal F$ is nonempty, choose $\varepsilon>0$ with
$3\varepsilon<\mathbf d(\mathcal R,\mathcal F)$; otherwise
choose any $\varepsilon>0$.  The fixed open set
\begin{equation}
 \mathcal O=\{\mu:\mathbf d(\mu,\mathcal R)<\varepsilon\}
 \label{tree:eq:bootstrap_success_neighborhood}
\end{equation}
contains $\mathcal R$, and every law in
$\overline{\mathcal O}\cap\mathcal P_{K_1}$ satisfies the
weaker shell tests.  Notice that $\mathcal O$ is open in the
ambient law space, not just relatively open in
$\mathcal P_{K_1}$.  This is needed when a failed scale is
renormalized only after its limiting norm has been found.

\paragraph{Successful scales.}
For $j\geq m$, say that level $j$ is \emph{successful} when
all three strict conditions hold:
\begin{equation}
 \mu_j\in\mathcal O,\qquad
 R^{-1}<\frac{A_j}{A_{j-1}}<R,\qquad
 \frac{A_{j-m}}{A_j}<2^{-ms}.
 \label{tree:eq:bootstrap_success}
\end{equation}
Success is a deterministic property of the law $\mu_j$ and the
normalizing sequence, rather than an event for a sampled map.
The topology, shell tests, $R,m$, and $\mathcal O$ are now fixed.
Proposition~\ref{tree:prop:internal_norm_comparison}
shows why these conditions are strict at every rigid limit,
including at any fixed forward or backward scale.  Indeed, at
such a scale its adjacent ratio is between
$R_0^{-1}2^p$ and $R_0 2^p$, while its $m$-step past ratio
is at most
\begin{equation}
 R_0 2^{-mp}<\tfrac12\,2^{-ms}.
 \label{tree:eq:bootstrap_strict_improvement}
\end{equation}

\subsection{No first failure after a late record}

At record scales, the complete past tail gives the moment bounds
needed for rigidity.  We first show that a late record starts an
arbitrarily long successful run.  If this run has a first failure,
its preceding scales still provide enough control to identify the
failed scale in the preceding normalization; it will then satisfy
all three success tests as well.

\begin{proposition}[Eventual control of all dyadic scales]
\label{tree:prop:bootstrap}
There are an integer $J$ and a finite deterministic $K$ such
that every $j\geq J$ satisfies \eqref{tree:eq:bootstrap_success},
and
\begin{equation}
 \frac{A_{j-i}}{A_j}\leq K2^{-si}
 \qquad(j\geq J,\ 0\leq i\leq j).
 \label{tree:eq:bootstrap_allscale_bound}
\end{equation}
In particular, $A_j\to\infty$.  Every sequence $j_r\to\infty$
has a subsequence whose joint normalized data converge to a
law in $\mathcal R$.
\end{proposition}

\begin{proof}
\smallskip\noindent\emph{Record scales start long successful runs.}
Lemma~\ref{tree:lem:record_scales} supplies $r_l\to\infty$ and
exponents $t_l\uparrow p$, with $s\leq t_l<p$, such that
\begin{equation}
 \frac{A_{r_l-i}}{A_{r_l}}\leq2^{-t_li}
 \qquad(0\leq i\leq r_l).
 \label{tree:eq:bootstrap_record_tail}
\end{equation}
These are record maxima of $2^{-t_lj}A_j$; their existence
uses only the volume-growth consequence
$\limsup j^{-1}\log_2 A_j\geq p$.
In any joint extraction at these records, each fixed-offset
internal norm is at most $2^{-ip}$.  The first alternative of
Proposition~\ref{tree:prop:comparison_families} and
Theorem~\ref{tree:thm:rigidity} apply.  The unit norm is one by
Lemma~\ref{tree:lem:moment_ui}, so the resulting law belongs to
$\mathcal R$.

Fix an integer $L$.  Extract simultaneously on the bounded
enlargement needed for levels $r_l-L,\ldots,r_l+L$ and their
$m$ preceding levels.  Initially keep cost units $A_{r_l}$
at all of these levels.  The moment lemmas give compactness
and moment convergence on this finite enlargement.  The
record limit is rigid, so
Proposition~\ref{tree:prop:internal_norm_comparison} gives a
strictly positive internal norm at every one of the enlarged
or reduced scales.  We may therefore divide by its own
limiting norm.  The resulting normalized laws are in
$\mathcal R$, and \eqref{tree:eq:bootstrap_constants}--\eqref{tree:eq:bootstrap_strict_improvement}
make their numerical success tests strict.  Their law tests
are strict since $\mathcal O$ is open and contains
$\mathcal R$.  If success failed infinitely often for this
fixed finite collection of levels, extracting that sequence
would contradict these conclusions.  Thus every fixed
finite run around a record eventually succeeds.  A diagonal
choice, relabelled $r_l$, makes all levels
$r_l,\ldots,r_l+l$ successful.  The complete record tail
\eqref{tree:eq:bootstrap_record_tail} is retained.

\smallskip\noindent\emph{The past tail along a successful run.}
Suppose $r\leq j$ and all levels from $r$ to $j$ are
successful, while $r$ has the record tail
$A_q/A_r\leq2^{-s(r-q)}$ for $q\leq r$.
For $r\leq q\leq j$, write $j-q=am+b$ with
$0\leq b<m$.  Iterating the $m$-step inequalities $a$
times and then the adjacent inequalities $b$ times gives
\[
 \frac{A_q}{A_j}\leq 2^{-ams}R^b
 \leq(R2^s)^{m-1}2^{-s(j-q)}.
\]
For $q<r$, apply the same estimate to $A_r/A_j$ and
multiply it by the record-tail estimate for $A_q/A_r$.
Thus in both cases
\begin{equation}
 \frac{A_q}{A_j}\leq K_0 2^{-s(j-q)}
 \qquad(0\leq q\leq j).
 \label{tree:eq:bootstrap_prefailure_bound}
\end{equation}
The constant has no dependence on the length of the
successful run.

\smallskip\noindent\emph{Extract immediately before a first failure.}
Suppose unsuccessful levels occur arbitrarily far out.
For each selected record $r_l$ let $f_l>r_l$
be the first unsuccessful level after it, and put
$j_l=f_l-1$.  The arbitrarily long initial successful
runs imply $j_l-r_l\to\infty$.
Use time units $2^{j_l}$ and cost units $A_{j_l}$.
Equation~\eqref{tree:eq:bootstrap_prefailure_bound} gives
compactness, the time modulus, and moment convergence.
For every fixed $i\geq0$, the normalized law at level
$j_l-i$ is successful for large $l$.  Its norm in the
base units has a positive finite subsequential limit:
the adjacent inequalities give
\[
 R^{-i}\leq A_{j_l-i}/A_{j_l}\leq R^i,
\]
and the sharper upper estimate is
\eqref{tree:eq:bootstrap_prefailure_bound}.  Take these limits
jointly for all fixed $i$.  Their laws lie in
$\overline{\mathcal O}$ and their moment bounds put them
in $\mathcal P_{K_1}$.  They therefore pass the coarse
shell tests.  The second alternative of
Proposition~\ref{tree:prop:comparison_families} applies to
this scale family.  It gives exact-power upper bounds
and compact zoom laws; Theorem~\ref{tree:thm:rigidity} then
makes the base law rigid.  Its unit internal norm is
one, and its deterministic comparison constant lies
in $[C_-,C_+]$.

\smallskip\noindent\emph{Recover the failed scale's normalization.}
Keep the preceding cost units $A_{j_l}$.  The unsuccessful level
is then the interval $[0,2]$, so it can be examined before any
lower bound on its own normalization is known.
Equation~\eqref{tree:eq:bootstrap_elementary_growth} gives
\begin{equation}
 0< A_{f_l}/A_{j_l}\leq2,
 \label{tree:eq:bootstrap_failed_upper_only}
\end{equation}
but no lower bound has been used.  Addition and
stationarity on this finite enlargement give the
same compactness and uniform integrability as before.
Thus, after extraction, its internal endpoint norm is
exactly
\[
 q_{-1}=\lim_l A_{f_l}/A_{j_l}.
\]
It is now positive by
Proposition~\ref{tree:prop:internal_norm_comparison}, applied
to the rigid law on the enlarged window:
\begin{equation}
 \alpha C2^p\leq q_{-1}\leq\beta C2^p.
 \label{tree:eq:bootstrap_failed_positive}
\end{equation}
Equivalently, one may transport the enlarged interval
to unit time and divide costs by $2^pA_{j_l}$; its
half-interval norm is then exactly $2^{-p}$, so
nondegeneracy is already present before the unit
interval norm is recovered.  Both descriptions use
preceding units first.

We may now normalize the failed-level data by
$A_{f_l}$.  Their limit is the time-$2$ transport of
the preceding rigid law divided by its positive
unit endpoint norm, hence belongs to $\mathcal R$.
Thus their law test passes eventually.  Furthermore,
moment convergence on the same enlarged window and
\eqref{tree:eq:bootstrap_norm_bridge} give
\[
 \lim_l\frac{A_{f_l}}{A_{f_l-1}}
 \in[R_0^{-1}2^p,R_0 2^p],\qquad
 \lim_l\frac{A_{f_l-m}}{A_{f_l}}
 \leq R_0 2^{-mp}.
\]
These limits lie strictly within the permitted
adjacent and $m$-step ranges.  All three tests at
$f_l$ therefore pass for large $l$, contradicting
its definition.  This proves eventual success.

\smallskip\noindent\emph{The global bound and all subsequences.}
Iterating the eventual $m$-step inequalities and using adjacent
ratios for the remainder gives
\eqref{tree:eq:bootstrap_allscale_bound} for $j-i\geq J$.
The finitely many values below $J$ are absorbed into
a single larger deterministic constant $K$ by
concatenating at $J$.  On each residue class modulo
$m$ we also have, for all sufficiently large $j$,
\[
 A_{j+m}>2^{ms}A_j.
\]
The finitely many residue classes have positive
starting values, proving $A_j\to\infty$.

Finally take any $j_r\to\infty$.
Equation~\eqref{tree:eq:bootstrap_allscale_bound} gives a
joint extraction with preserved unit norm.  At every
fixed past offset the success tests hold eventually.
The limiting fixed-offset bounds use $K_0$, since
their entire finite runs lie beyond $J$.  Thus the
closure property of $\mathcal O$ again supplies the
second alternative of
Proposition~\ref{tree:prop:comparison_families}.  Rigidity
and \eqref{tree:eq:bootstrap_C_range} put the resulting law
in $\mathcal R$, as claimed.
\end{proof}

\subsection{Arbitrary time units and the interval estimate}

The dyadic conclusion now extends to any integer time unit.
The remaining input is the already proved time modulus, which
allows a converging multiplicative change of time to be replaced
by its limit.  We also record the uniform interval estimate needed
for the finite-sphere argument.

\begin{theorem}[Bilateral subsequential limit at all time scales]
\label{tree:thm:bilateral_limit}
For a positive integer $N$ set
\[
 j(N)=\lfloor\log_2 N\rfloor,\qquad B_N=A_{j(N)}.
\]
Then $B_N\to\infty$.  From every sequence $N_r\to\infty$
one can extract a subsequence such that the bilateral
contours and all their local certificate data, in time
units $N_r$ and cost units $B_{N_r}$, converge jointly.
The continuum background is the fixed-unit
$\sqrt2$-quantum cone decorated by its area-parametrized
$\SLE_8$.  For one deterministic
\begin{equation}
 \lambda\in[C_-,2^p C_+],
 \label{tree:eq:bootstrap_lambda_range}
\end{equation}
its limiting local passage rule is
\begin{equation}
 F_U(x,y)=\lambda D_h(x,y;U)
 \label{tree:eq:bootstrap_bilateral_identification}
\end{equation}
for every open $U$ and every $x,y$ joined by a path
compactly in $U$.  The equality uses the internal
length metric on $U$, equivalently exhaustion by
compactly confined paths, and includes the upper
realization and lower subpath-extraction assertions
with open margins.

There is also the following uniform discrete estimate.
For each fixed $T<\infty$ there is $K_T<\infty$ such
that, for every integer $N\geq1$ and every translated
integer interval $I$ of length $|I|\leq TN$,
\begin{equation}
 \left\|\frac{\diam(V(G_I),d_I)}{B_N}\right\|_k
 \leq K_T\left(\frac{|I|+1}{N}\right)^s.
 \label{tree:eq:bootstrap_uniform_interval_diameter}
\end{equation}
The corresponding curve-time modulus on each bounded
window is uniform, with any exponent smaller than
$s-1/k$.  The fixed-interval endpoint $k$th moments
are uniformly integrable along these extractions.
\end{theorem}

\begin{proof}
\smallskip\noindent\emph{Identification in arbitrary time units.}
Divergence and the dyadic conclusion follow from
Proposition~\ref{tree:prop:bootstrap}. For arbitrary $N_r$, put
$j_r=j(N_r)$ and pass to a subsequence for which
\[
 u_r=N_r/2^{j_r}\longrightarrow u\in[1,2].
\]
Extract the dyadic laws jointly on every bounded enlargement.
Their limiting local rule is $C D_h$ for a deterministic
$C\in[C_-,C_+]$.

Proposition~\ref{tree:prop:zoom_transport} gives transport for
every fixed positive time factor. Having fixed the deterministic
$u$, adjoin to the retained countable family the $u$-transforms
of its templates and all their fixed dyadic offsets, and take a
further joint extraction. Retain also
\[
 (\mathcal C_{j_r},T_u\mathcal C_{j_r},r_u(\mathcal C_{j_r})),
\]
including the transformed internal reference metrics. Here
$\mathcal C_j$ is the coherent background of
Section~\ref{tree:sec:coherent-background}, and $r_u$ is the
circle-average dilation in
\eqref{tree:eq:background_joint_zoom}. This triple has the same
joint marginal for every $r$. Its limiting components therefore
have that same marginal and satisfy the same measurable transport
identity. No continuity of the chosen reembedding is required.

The raw paths in $N_r$ time units are the same paths with their
dyadic time labels divided by $u_r$; their costs remain in
$A_{j_r}=B_{N_r}$ units. These time changes converge uniformly on
bounded windows to division by $u$. For internal endpoint
distances, the uniform time modulus and
Lemma~\ref{tree:lem:cut_stability} identify the moving endpoints
$u_rt$ with $ut$, since the fixed limiting cuts have singleton
Brownian fibers. The same identification for local arrays follows
from finite-band reconstruction with open margins and both
directions of marked-certificate convergence.

To keep the spatial background in this comparison explicit,
let $\eta_j$ be the curve component of $\mathcal C_j$ and write
$r_{u,j}=r_u(\mathcal C_j)$ and
$\eta_j^{(u)}(t)=r_{u,j}^{-1}\eta_j(ut)$. In the fixed-$u$
spatial coordinates, an $N_r$ time label $t$ is projected to
\[
 r_{u,j_r}^{-1}\eta_{j_r}(u_rt)
       =\eta_{j_r}^{(u)}((u_r/u)t).
\]
Joint local uniform convergence of the retained curves makes
its difference from $\eta_{j_r}^{(u)}(t)$ tend uniformly to zero
on bounded time windows, in probability. Thus the limiting $N_r$ arrays are the
dyadic limiting arrays transported to the fixed-$u$ background,
with the original costs. The fixed conditional background law
is retained as in \eqref{tree:eq:bootstrap_background_closure}.

Write $r_u$ for the limiting dilation and $h^{(u)}$ for the
transformed field. The local rule in these coordinates is
therefore
\[
 F^{\mathrm{new}}_U(x,y)
   =C D_h(r_ux,r_uy;r_uU)
   =Cu^pD_{h^{(u)}}(x,y;U).
\]
The transformed background has the fixed unit cone law by
\eqref{tree:eq:background_joint_zoom}. Consequently
$\lambda=Cu^p$ is deterministic and satisfies
\eqref{tree:eq:bootstrap_lambda_range}. The factor $u^p$ comes
from the established reference metric's quantum-unit rule;
the passage law is transported by
Proposition~\ref{tree:prop:zoom_transport}.
Equality on all open sets follows first on the countable
buffered charts, then for compactly confined paths by finite
covering and concatenation, and finally by exhaustion.
These operations preserve upper realization and lower marked
subpath extraction with open margins, as asserted.

\smallskip\noindent\emph{Uniform interval bounds.}
Apply Lemma~\ref{tree:lem:moment_compactness} to
\eqref{tree:eq:bootstrap_allscale_bound}.  It gives,
uniformly in translated intervals of length at
most $T2^j$, the internal diameter bound
$K'_T((|I|+1)/2^j)^s$ in units $A_j$.
Since $N/2\leq2^{j(N)}\leq N$, enlarging the
constant gives
\eqref{tree:eq:bootstrap_uniform_interval_diameter}.
Finitely many small $N$ are absorbed into that
constant.  The same lemma gives the positive
time-modulus exponent, and
Lemma~\ref{tree:lem:moment_ui} gives uniform
integrability, also on the finite enlargements
used here.  Neither the argument nor the estimate
requires regular variation of $A_j$ or convergence
of its adjacent ratios.
\end{proof}
